\par\textbf{Procedimiento completo para comparar tu solución.} La gravedad $\mu/r^2$ proporciona aceleración centrípeta $v^2/r$. Despejar $v=\sqrt{\mu/r}$ y usar longitud orbital $2\pi r$ para periodo. La ingravidez aparente corresponde a caída libre, no gravedad nula.

\begin{enumerate}

\item Sumar radio terrestre y altura antes de convertir km.

\[\fitmath{r=\left(6371+500\right)\cdot 1000\approx 6871000\,\mathrm{m}}\]

\item Velocidad circular.

\[\fitmath{v=\sqrt{\frac{3.986004418\times10^{14}}{r}}}\]

\[\fitmath{v=\sqrt{\frac{3.986004418\times10^{14}}{(6871000)}}\approx 7616.561\,\mathrm{m/s}}\]

\item Tiempo de una vuelta.

\[\fitmath{P=\frac{2\cdot \pi\cdot r}{v}}\]

\[\fitmath{P=\frac{2\cdot \pi\cdot (6871000)}{(7616.560806)}\approx 5668.144\,\mathrm{s}}\]

\item Gravedad local.

\[\fitmath{g=\frac{3.986004418\times10^{14}}{{r}^{2}}}\]

\[\fitmath{g=\frac{3.986004418\times10^{14}}{{(6871000)}^{2}}\approx 8.443021\,\mathrm{m/s^{2}}}\]

\item Fuerza gravitatoria sobre el satélite.

\[\fitmath{F=100\cdot g}\]

\[\fitmath{F=100\cdot (8.443021178)\approx 844.3021\,\mathrm{N}}\]

\end{enumerate}
